% ID: SEL-D0027-OB-E01
% ESTADO: validada. S028, cotejada con examen y criterios.
\textbf{a)} En el primer cuadrante, $y$ queda por debajo de
$(11-x)/2$, $(x+5)/2$ y $18-3x$. Las fronteras que forman el
techo se cortan en $(3,4)$ y $(5,3)$, al igualar la segunda con
la primera y la primera con la tercera, respectivamente.
Los cortes factibles con los ejes son $(0,5/2)$ y $(6,0)$,
además del origen.
\begin{center}
\setlength{\unitlength}{8mm}
\begin{picture}(10,6)(-2,-1)
\qbezier(0,0)(3.5,0)(7,0)\put(7.2,0){$x$}
\qbezier(0,0)(0,2.25)(0,4.5)\put(0,4.6){$y$}
\qbezier(0,2.5)(1.5,3.25)(3,4)
\qbezier(3,4)(4,3.5)(5,3)
\qbezier(5,3)(5.5,1.5)(6,0)
\put(0,2.5){\makebox(0,0)[r]{$(0,5/2)$}}
\put(3,4.2){\makebox(0,0)[b]{$(3,4)$}}
\put(5.1,3.1){$(5,3)$}
\put(6,-0.5){\makebox(0,0)[t]{$(6,0)$}}
\put(0,-0.5){\makebox(0,0)[t]{$(0,0)$}}
\end{picture}
\end{center}
\textbf{b)} En $(5.5,2)$ se obtiene $3(5.5)+2=18.5>18$.
Basta esa desigualdad incumplida para excluir el punto.

\textbf{c)} En $(0,0)$, $(6,0)$, $(5,3)$, $(3,4)$ y $(0,5/2)$,
la función lineal vale $0,12,19,18,7.5$, respectivamente.
El máximo es $\boxed{19}$ sólo en $(5,3)$ y el mínimo
$\boxed{0}$ sólo en $(0,0)$.
